%% Sect. 8 -- from augur_AMT_§8 draft v1 (2026-09-22); rewritten for length and sentence length 2026-09-29.
\conclusions
\label{sec:8}

We have presented Augur, an open DOAS retrieval posed in the extinction domain as a separable nonlinear least-squares
problem and solved by variable projection. It is configured by instrument and campaign profiles, and its diagnostics
were demonstrated on one field case: a thermal-dissociation broadband cavity-enhanced absorption spectrometer with two
heated and one ambient-temperature cavity. The formulation is conventional; the implementation makes three things cheap: substituting a calibration reference
without reprocessing the campaign, evaluating the projected objective anywhere in the nonlinear parameters, and
carrying the termination state and the distance to each reference into the product.

The code solves the problem it states: synthetic recovery is at the level of the linear algebra, the analytic
Jacobian agrees with finite differences to round-off, and Augur and QDOAS agree on synthetic spectra. On field spectra they agree in NO$_2$ with
$r^2 \ge \fieldnum{0.990}$ but differ in the heated-channel scale by \fieldnum{4}--\fieldnum{6}\,\%. The full chain returns \fieldnum{$+0.0054$}\,ppb at a true zero, a
twentieth of the three-sigma detection limit of \fieldnum{0.203} to about \fieldnum{0.257}\,ppb, which tests the additive chain only.

The first question concerned errors the fit cannot see. With white noise and a complete model the reported sigma tracks
the synthetic scatter to within about ten per cent, rising to 1.32 with a free shift; against the calibration chain it
fails. For the channel-difference product, perturbing the
zero-air reference, the reflectivity and the etalon as the operational pipeline builds them gives, in the
best-calibrated week, a structural contribution of \fieldnum{1.28} (robust scale) to \fieldnum{2.29} (standard deviation) times the
reported sigma. The total uncertainty is \fieldnum{1.63} to \fieldnum{2.50} times what the fit reports for the same 60\,s records. None of the
calibration-chain terms enlarges the residual. These terms add to the campaign-constant cross-section and
Rayleigh uncertainties of earlier treatments rather than replacing them (Sect.~\ref{sec:4.6}). The magnitudes belong to
this instrument; the mechanism, an intermittently measured reference interpolated to every record, is general.

The second question concerned information the pipeline holds. With the cold shift bounded to $\pm 1$\,px in a
sensitivity run, every base fit of the ambient-temperature channel returned its shift at the bound. In the budget
propagation \fieldnum{187} of \fieldnum{9034} records were excluded because a fit terminated at a bound, and all were recorded as
converged. In one window the projected objective varies by
about one per cent across four pixels of shift, in another over five orders of magnitude, and both fits report
convergence. A reflectivity reference on a superseded time convention displaced the alkyl nitrates by \fieldnum{1.44} times the
structural budget, while the residual and the reported uncertainty moved by two parts in a thousand.

A third category lies outside both. Against a co-located NIER analyser the heated-channel NO$_2$ reads low by a factor of
about \fieldnum{0.6} while the ambient-temperature channel agrees. This multiplicative discrepancy, whose cause is unresolved, is
the largest error in the record, and neither the fit report nor the proposed fields see it (Sects.~\ref{sec:3.6} and
\ref{sec:6.5}).

We therefore propose that a retrieval report three numbers: the fit uncertainty, the structural term for the record's own
calibration context, and a context flag (Sect.~\ref{sec:6.1}). Every implementation already computes the flag's fields,
except the build time and time convention of a reference, which must be recorded when the reference is built.
Quality gates should be reported with the coverage they cost, because they remove coverage as well as variance
(Sect.~\ref{sec:6.4}). Where the interpolation error is set by knot noise, as in the 180\,$^{\circ}$C channel, averaging more scans per
calibration block is the lever to test; the within-block evidence is unfavourable and the decisive measurement has not
been made (Sect.~\ref{sec:6.2}).

The evidence is one instrument over one campaign. The code is not tied to that instrument (Appendix~\ref{app:A}), but
its behaviour on other configurations is not shown here. The reporting convention and the synthetic benchmark of 261
cases, which contains no retrieval code or field data, should carry beyond it; extending the benchmark across line
shapes, windows and pixel grids would separate what is general from what is ours.

Three extensions follow from the results. Fitting one shift jointly across the spectra of a calibration block
could remove the indeterminacy diagnosed here, given an absorber in the block (Sect.~\ref{sec:5.1}). Treating the zero-air
reference and the reflectivity as measured quantities with their own covariance would propagate the structural term
inside the fit (Sect.~\ref{sec:4.6}). And the raw light level of each spectrum, the proximate cause of all six ambient-temperature excursions, should travel
with the product together with the helium and zero-air intensities of each calibration block (Sect.~\ref{sec:6.5}).
